% Modèle de mémoire / thèse -- projet de démonstration
% Classe KOMA-Script scrbook, babel (français + anglais), biblatex + biber, pgfplots, TikZ.
% Le contenu mathématique est classique ; noms, université et jury sont des exemples à remplacer.
% Compilation : latexmk -pdf main.tex
\documentclass[a4paper,11pt,twoside,openany,BCOR=7mm,DIV=12,headings=small,numbers=noenddot,parskip=false]{scrbook}

\usepackage[T1]{fontenc}
\usepackage[utf8]{inputenc}
\usepackage[english,main=french]{babel}
\usepackage{csquotes}
\usepackage{amsmath,amssymb,amsthm,mathtools}
\usepackage{newpxtext,newpxmath}
\usepackage{microtype}
\usepackage{xcolor}
\usepackage{graphicx}
\usepackage{booktabs}
\usepackage{siunitx}
\sisetup{locale=FR}
\usepackage{enumitem}
\usepackage{tikz}
\usetikzlibrary{arrows.meta,calc,decorations.pathreplacing}
\usepackage{pgfplots}
\pgfplotsset{compat=1.18}
\usepackage{listings}
\usepackage[style=alphabetic,backend=biber,giveninits=true,doi=false,isbn=false,url=false]{biblatex}
\addbibresource{refs.bib}
\usepackage{caption}
\usepackage{subcaption}
\usepackage{scrlayer-scrpage}
\usepackage[colorlinks=true,linkcolor=wine,citecolor=wine,urlcolor=wine,pdftitle={Méthodes numériques pour les équations différentielles ordinaires},pdfauthor={Prénom Nom}]{hyperref}
\usepackage[nameinlink]{cleveref}

% ------------------------------------------------------------------ couleurs
\definecolor{wine}{HTML}{7F1D1D}
\definecolor{ink}{HTML}{1C1917}
\definecolor{sea}{HTML}{1D4E89}
\definecolor{moss}{HTML}{2F6B3A}
\definecolor{ochre}{HTML}{B7791F}
\definecolor{sand}{HTML}{F3ECDA}
\definecolor{line}{HTML}{DDD2B8}

% ------------------------------------------------------------------ titres
\renewcommand*{\chapterformat}{\mbox{\color{wine}\fontsize{46}{46}\selectfont\thechapter}\enskip{\color{line}\rule[-0.5ex]{0.9pt}{3.4ex}}\enskip}
\setkomafont{chapter}{\normalfont\huge\bfseries}
\setkomafont{section}{\normalfont\Large\bfseries}
\setkomafont{subsection}{\normalfont\large\bfseries}
\setkomafont{descriptionlabel}{\normalfont\bfseries}
\addtokomafont{chapterprefix}{\color{wine}}
\setkomafont{captionlabel}{\bfseries\color{wine}}
\captionsetup{font=small,labelsep=period}

% ------------------------------------------------------------------ en-têtes
\automark[section]{chapter}
\clearpairofpagestyles
\lehead{\pagemark\quad\headmark}
\rohead{\headmark\quad\pagemark}
\cfoot*{\pagemark}
\KOMAoptions{headsepline=true}
\setkomafont{pageheadfoot}{\normalfont\small\itshape}
\setkomafont{pagenumber}{\normalfont\small\bfseries\upshape}
\setkomafont{headsepline}{\color{line}}
\renewcommand*{\chapterpagestyle}{plain.scrheadings}
\pagestyle{scrheadings}

% ------------------------------------------------------------------ théorèmes
\theoremstyle{plain}
\newtheorem{theoreme}{Théorème}[chapter]
\newtheorem{proposition}[theoreme]{Proposition}
\newtheorem{corollaire}[theoreme]{Corollaire}
\theoremstyle{definition}
\newtheorem{definition}[theoreme]{Définition}
\newtheorem{exemple}[theoreme]{Exemple}
\theoremstyle{remark}
\newtheorem{remarque}[theoreme]{Remarque}
\crefname{theoreme}{théorème}{théorèmes}
\crefname{proposition}{proposition}{propositions}
\crefname{corollaire}{corollaire}{corollaires}
\crefname{definition}{définition}{définitions}
\crefname{exemple}{exemple}{exemples}
\crefname{remarque}{remarque}{remarques}
\crefname{chapter}{chapitre}{chapitres}
\crefname{section}{section}{sections}
\crefname{figure}{figure}{figures}
\crefname{table}{tableau}{tableaux}
\crefname{appendix}{annexe}{annexes}

% ------------------------------------------------------------------ raccourcis
\newcommand{\R}{\mathbb{R}}
\newcommand{\N}{\mathbb{N}}
\newcommand{\norm}[1]{\lVert #1\rVert}
\newcommand{\dd}{\mathrm{d}}

% ------------------------------------------------------------------ code source
\lstset{
  basicstyle=\ttfamily\small, keywordstyle=\bfseries\color{wine}, commentstyle=\itshape\color{ink!60},
  stringstyle=\color{sea}, language=Python, frame=single, rulecolor=\color{line},
  backgroundcolor=\color{sand!50}, numbers=left, numberstyle=\tiny\color{ink!50}, xleftmargin=1.6em,
  showstringspaces=false, columns=fullflexible, keepspaces=true,
}

% ------------------------------------------------------------------ style des graphiques
\pgfplotsset{
  these/.style={
    width=0.92\linewidth, height=6.2cm, font=\footnotesize,
    tick label style={font=\scriptsize}, label style={font=\footnotesize},
    axis line style={ink!70}, tick style={ink!70},
    grid=major, grid style={line!70, line width=0.3pt},
    legend style={font=\scriptsize, fill=white, fill opacity=0.9, draw=line, text opacity=1, cells={anchor=west}},
    every axis plot/.append style={thick},
  },
}

\setlength{\emergencystretch}{2em}

\begin{document}

\frontmatter
\pagestyle{empty}
\begin{titlepage}
  \centering
  {\scshape\Large Université Exemple}\\[3pt]
  {\small Faculté des Sciences et Techniques\\ Département de Mathématiques}\par
  \vspace{8pt}
  {\color{wine}\rule{\textwidth}{1.2pt}}\par
  \vspace{2.4cm}

  {\small\scshape Mémoire de Master en Mathématiques appliquées}\par
  \vspace{1.1cm}
  {\fontsize{26}{32}\selectfont\bfseries Méthodes numériques pour les équations différentielles ordinaires\par}
  \vspace{10pt}
  {\Large\itshape Analyse de convergence et applications\par}

  \vspace{1.6cm}
  \begin{tikzpicture}
    \draw[line,line width=0.8pt] (0,0) -- (3.2,0);
    \fill[wine] (1.6,0) circle (2.2pt);
  \end{tikzpicture}\par

  \vspace{1.6cm}
  {\small Présenté par}\par
  \vspace{4pt}
  {\LARGE Prénom Nom\par}

  \vspace{1.6cm}
  {\small Sous la direction de}\par
  \vspace{4pt}
  {\large Pr.~Nom du directeur\par}
  {\small Université Exemple}\par

  \vfill
  \begin{tabular}{@{}rl@{}}
    \itshape Président du jury : & Pr.~Nom Un\\
    \itshape Rapporteur :        & Dr~Nom Deux\\
    \itshape Examinateur :       & Dr~Nom Trois
  \end{tabular}\par
  \vspace{1.2cm}
  {\small Octobre 2026}
\end{titlepage}

\cleardoublepage
\pagestyle{plain.scrheadings}
\chapter*{Remerciements}
\addcontentsline{toc}{chapter}{Remerciements}

Je tiens à remercier mon directeur de mémoire, le professeur Nom du directeur, pour ses conseils, sa disponibilité et la confiance qu'il m'a accordée tout au long de ce travail.

Mes remerciements vont également aux membres du jury, qui ont accepté d'évaluer ce mémoire, ainsi qu'à l'ensemble de l'équipe pédagogique du département de mathématiques.

Enfin, je remercie ma famille et mes proches pour leur soutien constant.

\medskip
\begin{flushright}
  \itshape Prénom Nom
\end{flushright}

\vfill
{\small\itshape Ce texte est un exemple : remplacez-le par vos propres remerciements.}

\chapter*{Résumé}
\addcontentsline{toc}{chapter}{Résumé}

Ce mémoire étudie les méthodes numériques à un pas pour la résolution des équations différentielles ordinaires. Après avoir rappelé le cadre théorique du problème de Cauchy (existence et unicité de la solution), nous présentons trois méthodes explicites classiques : la méthode d'Euler, la méthode de Heun et la méthode de Runge--Kutta d'ordre~4.

Nous démontrons qu'une méthode consistante d'ordre~$p$, dont la fonction d'incrément est lipschitzienne, converge à l'ordre~$p$, et nous donnons une estimation explicite de l'erreur globale. Ces résultats sont ensuite confrontés à des expériences numériques : sur une équation test, les ordres observés valent $1$, $2$ et $4$, conformément à la théorie, et le coût nécessaire pour atteindre une erreur de $10^{-6}$ varie de plus de quatre ordres de grandeur d'une méthode à l'autre. Enfin, nous appliquons la méthode de Runge--Kutta au modèle proie-prédateur de Lotka--Volterra et vérifions la conservation de son intégrale première.

\bigskip
\noindent\textbf{Mots-clés :} équations différentielles ordinaires, méthodes à un pas, Runge--Kutta, ordre de convergence, erreur globale.

\clearpage
\begin{otherlanguage}{english}
\chapter*{Abstract}
\addcontentsline{toc}{chapter}{Abstract}

This thesis studies one-step numerical methods for ordinary differential equations. After recalling the theoretical framework of the Cauchy problem (existence and uniqueness of the solution), we present three classical explicit methods: Euler's method, Heun's method and the fourth-order Runge--Kutta method.

We prove that a consistent method of order~$p$ whose increment function is Lipschitz continuous converges with order~$p$, and we give an explicit bound on the global error. These results are then compared with numerical experiments: on a test equation the observed orders are $1$, $2$ and $4$, as predicted, and the cost required to reach an error of $10^{-6}$ differs by more than four orders of magnitude between the methods. Finally, we apply the Runge--Kutta method to the Lotka--Volterra predator--prey model and check the conservation of its first integral.

\bigskip
\noindent\textbf{Keywords:} ordinary differential equations, one-step methods, Runge--Kutta, order of convergence, global error.
\end{otherlanguage}


\tableofcontents
\listoffigures
\listoftables

\mainmatter
\pagestyle{scrheadings}
\chapter*{Introduction générale}
\addcontentsline{toc}{chapter}{Introduction générale}
\markboth{Introduction générale}{Introduction générale}

Les équations différentielles ordinaires (EDO) interviennent dans la modélisation de nombreux phénomènes : dynamique des populations, mécanique, circuits électriques, cinétique chimique. Dans la grande majorité des cas, leur solution ne s'exprime pas à l'aide de fonctions usuelles, et l'on doit recourir à des méthodes numériques pour l'approcher.

Les premières méthodes de ce type remontent à Euler~\cite{euler1768}. Elles ont été considérablement améliorées par Runge, Heun et Kutta~\cite{runge1895,heun1900,kutta1901}, dont les schémas restent aujourd'hui au c\oe ur de la plupart des logiciels de calcul scientifique. La théorie moderne de leur convergence est exposée dans plusieurs ouvrages de référence~\cite{hairer1993,butcher2016,suli2003,atkinson2009}.

\paragraph{Objectifs.}
Ce mémoire poursuit trois objectifs :
\begin{enumerate}[label=(\roman*),itemsep=2pt]
  \item rappeler le cadre théorique du problème de Cauchy et illustrer son intérêt sur le modèle proie-prédateur de Lotka et Volterra~\cite{lotka1925,volterra1926} ;
  \item présenter trois méthodes à un pas (Euler, Heun, Runge--Kutta d'ordre~4) et démontrer un théorème de convergence avec estimation de l'erreur globale ;
  \item confronter la théorie à des expériences numériques : ordres de convergence observés, coût de calcul, comportement sur un système non linéaire.
\end{enumerate}

\paragraph{Plan du mémoire.}
Le \cref{ch-edo} rappelle le problème de Cauchy, le théorème de Cauchy--Lipschitz et le modèle proie-prédateur. Le \cref{ch-methodes} introduit les méthodes à un pas, les notions de consistance et d'ordre, démontre le théorème de convergence et discute la stabilité absolue. Le \cref{ch-exp} présente les résultats numériques. Une conclusion propose quelques perspectives. Les tableaux de Butcher et un exemple de code figurent en annexe.

\paragraph{Données et reproductibilité.}
Toutes les données numériques de ce mémoire sont calculées et non mesurées. Les fichiers \texttt{convergence.csv} et \texttt{predator\_prey.csv}, fournis avec les sources, sont lus directement par \LaTeX{} pour tracer les graphiques.

\chapter{Équations différentielles ordinaires}\label{ch-edo}

Ce chapitre fixe le cadre et les notations utilisés dans la suite. Nous rappelons le problème de Cauchy, le théorème d'existence et d'unicité de Cauchy--Lipschitz, puis nous présentons un modèle non linéaire, celui de Lotka--Volterra, qui servira de fil conducteur.

\section{Le problème de Cauchy}\label{sec-cauchy}

Soient $d\ge 1$ un entier, $[t_0,T]$ un intervalle de $\R$ et $f\colon [t_0,T]\times\R^d\to\R^d$ une fonction continue. Le \emph{problème de Cauchy} consiste à trouver une fonction dérivable $y\colon[t_0,T]\to\R^d$ telle que
\begin{equation}\label{eq-cauchy}
  y'(t)=f\bigl(t,y(t)\bigr)\quad\text{pour tout } t\in[t_0,T],\qquad y(t_0)=y_0,
\end{equation}
où $y_0\in\R^d$ est la donnée initiale.

\begin{exemple}[Croissance exponentielle]\label[exemple]{ex-exp}
Pour $d=1$ et $f(t,y)=\lambda y$ avec $\lambda\in\R$, le problème \eqref{eq-cauchy} a pour solution
\[
  y(t)=y_0\,e^{\lambda(t-t_0)} .
\]
Nous utiliserons le cas $\lambda=1$, $y_0=1$, $t_0=0$ et $T=1$ comme problème test au \cref{ch-exp} : la solution exacte $y(t)=e^{t}$ permet de mesurer précisément l'erreur des méthodes.
\end{exemple}

\section{Existence et unicité}\label{sec-cl}

\begin{definition}[Fonction lipschitzienne]\label[definition]{def-lipschitz}
On dit que $f$ est \emph{lipschitzienne par rapport à $y$}, de constante $L\ge0$, si
\begin{equation}\label{eq-lipschitz}
  \norm{f(t,u)-f(t,v)}\le L\,\norm{u-v}\qquad\text{pour tous } t\in[t_0,T],\; u,v\in\R^d .
\end{equation}
\end{definition}

\begin{theoreme}[Cauchy--Lipschitz]\label{thm-cl}
Si $f$ est continue et lipschitzienne par rapport à $y$, alors le problème \eqref{eq-cauchy} admet une unique solution définie sur tout l'intervalle $[t_0,T]$.
\end{theoreme}

\begin{proof}[Idée de la démonstration]
Une fonction continue $y$ est solution de \eqref{eq-cauchy} si et seulement si elle est point fixe de l'application
\[
  (\mathcal{T}y)(t)=y_0+\int_{t_0}^{t} f\bigl(s,y(s)\bigr)\,\dd s
\]
sur l'espace $E=C([t_0,T],\R^d)$. Pour $\theta>L$, munissons $E$ de la norme $\norm{y}_\theta=\max_t e^{-\theta(t-t_0)}\norm{y(t)}$, qui est équivalente à la norme uniforme. Pour $y,z$ continues, le caractère lipschitzien donne
\[
  e^{-\theta(t-t_0)}\norm{(\mathcal{T}y)(t)-(\mathcal{T}z)(t)}
  \le L\,\norm{y-z}_\theta\, e^{-\theta(t-t_0)}\int_{t_0}^{t}e^{\theta(s-t_0)}\,\dd s
  \le \frac{L}{\theta}\,\norm{y-z}_\theta ,
\]
donc $\mathcal{T}$ est contractante et le théorème du point fixe de Banach conclut.
\end{proof}

Dans toute la suite, nous supposons que les hypothèses du \cref{thm-cl} sont satisfaites : le problème \eqref{eq-cauchy} est alors bien posé et la quantité $y(t_n)$ que l'on cherche à approcher est bien définie.

\section{Un modèle non linéaire : proie-prédateur}\label{sec-lv}

Le modèle de Lotka--Volterra~\cite{lotka1925,volterra1926} décrit l'évolution conjointe d'une population de proies $x(t)$ et d'une population de prédateurs $y(t)$ :
\begin{equation}\label{eq-lv}
  \begin{cases}
    x'=\alpha x-\beta xy,\\[2pt]
    y'=\delta xy-\gamma y,
  \end{cases}
  \qquad \alpha,\beta,\gamma,\delta>0 .
\end{equation}
Les proies se reproduisent à un taux $\alpha$ et sont consommées proportionnellement au produit $xy$ ; les prédateurs disparaissent à un taux $\gamma$ et se développent grâce aux proies consommées.

Le système \eqref{eq-lv} possède deux points d'équilibre : $(0,0)$ et
\[
  (x^\ast,y^\ast)=\Bigl(\frac{\gamma}{\delta},\,\frac{\alpha}{\beta}\Bigr).
\]

\begin{proposition}[Intégrale première]\label[proposition]{prop-invariant}
Pour toute solution de \eqref{eq-lv} avec $x,y>0$, la quantité
\[
  I(x,y)=\delta x-\gamma\ln x+\beta y-\alpha\ln y
\]
est constante au cours du temps.
\end{proposition}

\begin{proof}
Par dérivation composée,
\[
  \frac{\dd}{\dd t}I(x,y)=\Bigl(\delta-\frac{\gamma}{x}\Bigr)x'+\Bigl(\beta-\frac{\alpha}{y}\Bigr)y'
  =(\delta x-\gamma)(\alpha-\beta y)+(\beta y-\alpha)(\delta x-\gamma)=0 .\qedhere
\]
\end{proof}

\begin{remarque}
Ce système n'admet pas de solution exprimable avec des fonctions usuelles, mais la \cref{prop-invariant} fournit un moyen de contrôle : les trajectoires sont des courbes fermées autour de $(x^\ast,y^\ast)$, et une méthode numérique de bonne qualité doit conserver $I$ à la précision attendue. Nous utiliserons cette propriété au \cref{sec-exp-lv}.
\end{remarque}

\chapter{Méthodes numériques à un pas}\label{ch-methodes}

Nous nous plaçons dans le cadre du \cref{ch-edo} : $f$ est continue et lipschitzienne par rapport à $y$, et $y$ désigne l'unique solution du problème de Cauchy \eqref{eq-cauchy}. Ce chapitre présente les méthodes à un pas, leur ordre, et le théorème de convergence qui relie l'ordre local à la précision globale.

\section{Principe général}\label{sec-principe}

Soit $N\in\N^\ast$. On pose $h=(T-t_0)/N$ et $t_n=t_0+nh$ pour $n=0,\dots,N$. Une \emph{méthode à un pas} calcule des approximations $y_n\approx y(t_n)$ par la récurrence
\begin{equation}\label{eq-onestep}
  y_{n+1}=y_n+h\,\Phi(t_n,y_n;h),\qquad y_0\ \text{donné},
\end{equation}
où $\Phi$ est la \emph{fonction d'incrément}. Le calcul de $y_{n+1}$ n'utilise que $y_n$, d'où le nom de la méthode.

Deux erreurs doivent être distinguées (\cref{fig-erreur-locale}) :
\begin{itemize}[itemsep=2pt]
  \item l'\emph{erreur locale}, commise en un pas à partir d'une valeur exacte : $\delta(t;h)=y(t+h)-y(t)-h\,\Phi\bigl(t,y(t);h\bigr)$ ;
  \item l'\emph{erreur globale}, accumulée sur $n$ pas : $e_n=y(t_n)-y_n$.
\end{itemize}

\begin{figure}[tb]
  \centering
  \begin{tikzpicture}[x=2.6cm,y=0.8cm,>={Stealth[length=2.2mm]}]
    % axes
    \draw[->,ink!60] (0,0) -- (3.6,0) node[right] {\small $t$};
    \draw[->,ink!60] (0,0) -- (0,7.2) node[above] {\small $y$};
    % solution exacte
    \draw[thick,sea,domain=0:3.4,samples=50,smooth,variable=\x] plot ({\x},{1+0.6*\x+0.35*\x*\x});
    \node[sea,right,font=\small] at (3.4,6.9) {$y(t)$};
    % tangente et pas d'Euler
    \draw[ochre,dashed,thick] (0.35,{1.95+1.3*(0.35-1)}) -- (1,1.95);
    \draw[wine,very thick] (1,1.95) -- (2.2,3.51);
    % repères
    \draw[ink!50,dotted] (1,0) -- (1,1.95);
    \draw[ink!50,dotted] (2.2,0) -- (2.2,3.51);
    \node[below,font=\small] at (1,0) {$t_n$};
    \node[below,font=\small] at (2.2,0) {$t_{n+1}$};
    \draw[<->,ink!70] (1,-0.9) -- node[below,font=\small] {$h$} (2.2,-0.9);
    % points
    \fill[sea] (1,1.95) circle (2pt) node[above left,font=\small] {$y(t_n)=y_n$};
    \fill[wine] (2.2,3.51) circle (2pt) node[below right,font=\small] {$y_{n+1}$};
    \fill[sea] (2.2,4.014) circle (2pt) node[above left,font=\small] {$y(t_{n+1})$};
    % erreur locale
    \draw[decorate,decoration={brace,amplitude=4pt,mirror},wine] (2.2,3.51) -- (2.2,4.014);
    \node[wine,font=\small,align=left,anchor=west] at (2.42,3.76) {erreur locale};
  \end{tikzpicture}
  \caption{Un pas de la méthode d'Euler à partir d'un point exact $y(t_n)$ : l'écart entre $y(t_{n+1})$ et $y_{n+1}$ est l'erreur locale $\delta(t_n;h)$.}\label{fig-erreur-locale}
\end{figure}

\section{Trois méthodes classiques}\label{sec-classiques}

\paragraph{Méthode d'Euler explicite.}
Le développement de Taylor $y(t+h)=y(t)+h\,y'(t)+O(h^2)$ et l'équation $y'=f(t,y)$ conduisent à $\Phi(t,y;h)=f(t,y)$, soit
\[
  y_{n+1}=y_n+h\,f(t_n,y_n).
\]

\paragraph{Méthode de Heun.}
Elle remplace la pente en $t_n$ par la moyenne des pentes aux deux extrémités du pas, la valeur en $t_{n+1}$ étant prédite par Euler :
\[
  \Phi(t,y;h)=\tfrac12\Bigl[f(t,y)+f\bigl(t+h,\;y+h\,f(t,y)\bigr)\Bigr].
\]

\paragraph{Méthode de Runge--Kutta d'ordre~4 (RK4).}
On calcule quatre pentes intermédiaires
\begin{align*}
  k_1&=f(t,y), & k_2&=f\bigl(t+\tfrac h2,\;y+\tfrac h2k_1\bigr),\\
  k_3&=f\bigl(t+\tfrac h2,\;y+\tfrac h2k_2\bigr), & k_4&=f\bigl(t+h,\;y+h\,k_3\bigr),
\end{align*}
puis $\Phi(t,y;h)=\tfrac16\,(k_1+2k_2+2k_3+k_4)$. Le tableau de Butcher de cette méthode figure en \cref{app-butcher}. Le \cref{tab-methodes} résume le coût par pas et l'ordre de chacune.

\begin{table}[tb]
  \centering
  \caption{Les trois méthodes : coût par pas et ordre.}\label{tab-methodes}
  \small
  \begin{tabular}{@{}lcc@{}}
    \toprule
    Méthode & Évaluations de $f$ par pas & Ordre $p$\\
    \midrule
    Euler explicite & 1 & 1\\
    Heun            & 2 & 2\\
    Runge--Kutta classique (RK4) & 4 & 4\\
    \bottomrule
  \end{tabular}
\end{table}

\section{Consistance et ordre}\label{sec-consistance}

\begin{definition}[Consistance d'ordre $p$]\label[definition]{def-consistance}
La méthode \eqref{eq-onestep} est \emph{consistante d'ordre $p$} s'il existe des constantes $C,h_0>0$ telles que, pour la solution exacte $y$,
\begin{equation}\label{eq-consistance}
  \norm{\delta(t;h)}=\norm{y(t+h)-y(t)-h\,\Phi\bigl(t,y(t);h\bigr)}\le C\,h^{p+1}
\end{equation}
pour tous $0<h\le h_0$ et $t\in[t_0,T-h]$.
\end{definition}

\begin{proposition}[Ordre de la méthode d'Euler]\label[proposition]{prop-ordre-euler}
Si $y$ est de classe $C^2$ et $M=\max_t\norm{y''(t)}$, la méthode d'Euler est consistante d'ordre~$1$ avec $C=M/2$.
\end{proposition}

\begin{proof}
Par la formule de Taylor avec reste intégral,
\[
  y(t+h)-y(t)-h\,y'(t)=\int_t^{t+h}(t+h-s)\,y''(s)\,\dd s,
\]
et la norme de cette quantité est majorée par $M\int_t^{t+h}(t+h-s)\,\dd s=\frac{M}{2}h^2$.
\end{proof}

Les ordres de Heun et de RK4, égaux à $2$ et $4$, s'obtiennent par un développement de Taylor à des ordres plus élevés ; nous renvoyons à~\cite{hairer1993,butcher2016} pour les détails, qui font intervenir les conditions d'ordre des méthodes de Runge--Kutta.

\section{Convergence}\label{sec-convergence}

\begin{theoreme}[Convergence et erreur globale]\label{thm-convergence}
Supposons la méthode \eqref{eq-onestep} consistante d'ordre~$p$ et sa fonction d'incrément lipschitzienne par rapport à $y$, de constante $\Lambda>0$, uniformément en $t$ et en $h\le h_0$. Alors, avec $y_0=y(t_0)$, l'erreur globale vérifie
\begin{equation}\label{eq-erreur-globale}
  \norm{e_n}\le\frac{C}{\Lambda}\Bigl(e^{\Lambda(t_n-t_0)}-1\Bigr)h^{p}\qquad(0\le n\le N,\ h\le h_0).
\end{equation}
En particulier $\max_n\norm{e_n}=O(h^p)$ : la méthode converge à l'ordre~$p$.
\end{theoreme}

\begin{proof}
Notons $\delta_n=\delta(t_n;h)$. L'identité $y(t_{n+1})=y(t_n)+h\,\Phi(t_n,y(t_n);h)+\delta_n$, soustraite à \eqref{eq-onestep}, donne
\[
  e_{n+1}=e_n+h\bigl[\Phi(t_n,y(t_n);h)-\Phi(t_n,y_n;h)\bigr]+\delta_n .
\]
Les hypothèses sur $\Phi$ et \eqref{eq-consistance} entraînent
\[
  \norm{e_{n+1}}\le(1+h\Lambda)\norm{e_n}+C\,h^{p+1}.
\]
Comme $e_0=0$, une récurrence immédiate fournit
\[
  \norm{e_n}\le C\,h^{p+1}\sum_{k=0}^{n-1}(1+h\Lambda)^k
  =\frac{C\,h^{p}}{\Lambda}\Bigl((1+h\Lambda)^n-1\Bigr).
\]
Enfin $(1+h\Lambda)^n\le e^{nh\Lambda}=e^{\Lambda(t_n-t_0)}$, d'où \eqref{eq-erreur-globale}.
\end{proof}

\begin{corollaire}[Méthode d'Euler]\label[corollaire]{cor-euler}
Si $y\in C^2$, la méthode d'Euler vérifie
\[
  \norm{e_n}\le\frac{M}{2L}\Bigl(e^{L(t_n-t_0)}-1\Bigr)h,
\]
où $L$ est la constante de Lipschitz de $f$ et $M=\max_t\norm{y''(t)}$.
\end{corollaire}

\begin{proof}
Ici $\Phi=f$, donc $\Lambda=L$ ; la \cref{prop-ordre-euler} donne $p=1$ et $C=M/2$. On applique le \cref{thm-convergence}.
\end{proof}

\begin{remarque}
Le facteur $e^{\Lambda(T-t_0)}$ traduit l'amplification maximale des perturbations sur l'intervalle ; la borne est donc pessimiste sur de longs intervalles de temps. L'exposant $p$, en revanche, est celui que l'on observe en pratique (\cref{sec-exp-test}).
\end{remarque}

\section{Stabilité absolue}\label{sec-stabilite}

Appliquée à l'équation test $y'=\lambda y$ avec $\lambda\in\mathbb{C}$, chacune des trois méthodes s'écrit $y_{n+1}=R(z)\,y_n$ avec $z=h\lambda$, où
\[
  R_{\mathrm{Euler}}(z)=1+z,\qquad
  R_{\mathrm{Heun}}(z)=1+z+\tfrac{z^2}{2},\qquad
  R_{\mathrm{RK4}}(z)=1+z+\tfrac{z^2}{2}+\tfrac{z^3}{6}+\tfrac{z^4}{24}.
\]
La solution numérique reste bornée si et seulement si $|R(z)|\le1$. Pour $\lambda<0$ réel, cette condition impose $h\le 2/|\lambda|$ pour Euler et Heun ; la limite est un peu plus grande pour RK4 (environ $2{,}78/|\lambda|$). C'est cette contrainte de stabilité, et non la précision, qui limite le pas pour les problèmes raides : nous y reviendrons dans les perspectives.

\chapter{Expériences numériques}\label{ch-exp}

Ce chapitre confronte les résultats théoriques du \cref{ch-methodes} à des calculs effectués en double précision. Toutes les valeurs proviennent de programmes écrits pour ce mémoire : ce sont des données calculées, non des mesures.

\section{Problème test}\label{sec-exp-test}

Nous considérons l'\cref{ex-exp} avec $\lambda=1$ : $y'=y$, $y(0)=1$ sur $[0,1]$, de solution exacte $y(t)=e^t$ et donc $y(1)=e$. La constante de Lipschitz vaut $L=1$. Pour chaque méthode, nous intégrons avec $h=2^{-k}$, $k=2,\dots,8$, et nous relevons l'erreur $|e-y_N|$ au temps final. L'\emph{ordre observé} entre deux pas consécutifs est
\[
  p_{\mathrm{obs}}(h)=\log_2\frac{e(2h)}{e(h)} ,
\]
où $e(h)$ désigne l'erreur obtenue avec le pas $h$.

\begin{table}[htbp]
  \centering
  \caption{Erreur au temps $t=1$ et ordre observé pour $y'=y$, $y(0)=1$.}\label{tab-erreurs}
  \footnotesize
  \begin{tabular}{@{}l
    S[table-format=1.3e-1] S[table-format=1.2]
    S[table-format=1.3e-1] S[table-format=1.2]
    S[table-format=1.3e-2] S[table-format=1.2]@{}}
    \toprule
    & \multicolumn{2}{c}{Euler} & \multicolumn{2}{c}{Heun} & \multicolumn{2}{c}{RK4}\\
    \cmidrule(lr){2-3}\cmidrule(lr){4-5}\cmidrule(l){6-7}
    {$h$} & {erreur} & {ordre} & {erreur} & {ordre} & {erreur} & {ordre}\\
    \midrule
    1/4   & 2.769e-1 & {--}  & 2.343e-2 & {--}  & 7.189e-5  & {--}  \\
    1/8   & 1.525e-1 & 0.86  & 6.441e-3 & 1.86  & 4.984e-6  & 3.85  \\
    1/16  & 8.035e-2 & 0.92  & 1.688e-3 & 1.93  & 3.281e-7  & 3.93  \\
    1/32  & 4.129e-2 & 0.96  & 4.322e-4 & 1.97  & 2.105e-8  & 3.96  \\
    1/64  & 2.094e-2 & 0.98  & 1.093e-4 & 1.98  & 1.333e-9  & 3.98  \\
    1/128 & 1.054e-2 & 0.99  & 2.749e-5 & 1.99  & 8.384e-11 & 3.99  \\
    1/256 & 5.290e-3 & 0.99  & 6.893e-6 & 2.00  & 5.261e-12 & 3.99  \\
    \bottomrule
  \end{tabular}
\end{table}

Le \cref{tab-erreurs} montre que les ordres observés tendent vers $1$, $2$ et $4$ lorsque $h$ diminue, en accord avec le \cref{thm-convergence}. La \cref{fig-convergence} présente les mêmes données en échelles logarithmiques : les pentes des droites sont les ordres des méthodes.

\begin{figure}[htbp]
  \centering
  \begin{tikzpicture}
  \begin{loglogaxis}[
    these, xlabel={pas $h$}, ylabel={erreur en $t=1$},
    xmin=3e-3, xmax=0.4, ymin=1e-14, ymax=1, legend pos=south east,
    xtick={0.004,0.0156,0.0625,0.25}, xticklabels={$2^{-8}$,$2^{-6}$,$2^{-4}$,$2^{-2}$},
    ytick={1e-12,1e-9,1e-6,1e-3,1},
  ]
    \addplot[wine,mark=square*,mark size=1.6pt] table[col sep=comma,x=h,y=euler] {convergence.csv};
    \addlegendentry{Euler ($p=1$)}
    \addplot[sea,mark=*,mark size=1.6pt] table[col sep=comma,x=h,y=heun] {convergence.csv};
    \addlegendentry{Heun ($p=2$)}
    \addplot[moss,mark=triangle*,mark size=2pt] table[col sep=comma,x=h,y=rk4] {convergence.csv};
    \addlegendentry{RK4 ($p=4$)}
  \end{loglogaxis}
  \end{tikzpicture}
  \caption{Erreur au temps $t=1$ en fonction du pas pour $y'=y$, $y(0)=1$.}\label{fig-convergence}
\end{figure}

\section{Coût de calcul}\label{sec-exp-cout}

La précision a un prix : Euler demande une évaluation de $f$ par pas, Heun deux et RK4 quatre (\cref{tab-methodes}). Pour atteindre une erreur inférieure à $10^{-6}$ au temps $t=1$, le nombre minimal de pas, déterminé par dichotomie sur la formule exacte de récurrence $y_N=R(h)^N$, est donné dans le \cref{tab-cout}.

\begin{table}[htbp]
  \centering
  \caption{Coût pour obtenir une erreur inférieure à $10^{-6}$ sur $y'=y$ en $t=1$.}\label{tab-cout}
  \small
  \begin{tabular}{@{}l S[table-format=7.0] S[table-format=7.0]@{}}
    \toprule
    Méthode & {Nombre de pas} & {Évaluations de $f$}\\
    \midrule
    Euler & 1359469 & 1359469 \\
    Heun  & 673     & 1346    \\
    RK4   & 13      & 52      \\
    \bottomrule
  \end{tabular}
\end{table}

À cette tolérance, RK4 est plus de quatre ordres de grandeur moins coûteuse qu'Euler ($1{,}36\times10^{6}$ évaluations contre $52$), ce qui justifie l'usage de méthodes d'ordre élevé dès que l'on vise une précision modeste.

\section{Le modèle proie-prédateur}\label{sec-exp-lv}

Nous intégrons le système \eqref{eq-lv} avec la méthode RK4, pour les paramètres
\[
  \alpha=1,\quad \beta=0{,}1,\quad \gamma=1{,}5,\quad \delta=0{,}075,
  \qquad (x_0,y_0)=(10,\,5),
\]
un pas $\Delta t=0{,}01$ et un temps final $T=30$. Le point d'équilibre est $(x^\ast,y^\ast)=(20,\,10)$.

\begin{figure}[htbp]
  \centering
  \begin{subfigure}[t]{0.49\textwidth}
    \centering
    \begin{tikzpicture}
    \begin{axis}[
      these, width=\linewidth, height=5.6cm,
      xlabel={temps $t$}, ylabel={population}, xmin=0, xmax=30, ymin=0, legend pos=north east,
    ]
      \addplot[sea,no marks] table[col sep=comma,x=t,y=prey] {predator_prey.csv};
      \addlegendentry{proies}
      \addplot[wine,no marks] table[col sep=comma,x=t,y=predator] {predator_prey.csv};
      \addlegendentry{prédateurs}
    \end{axis}
    \end{tikzpicture}
    \caption{Évolution des populations}
  \end{subfigure}\hfill
  \begin{subfigure}[t]{0.49\textwidth}
    \centering
    \begin{tikzpicture}
    \begin{axis}[
      these, width=\linewidth, height=5.6cm,
      xlabel={proies}, ylabel={prédateurs}, xmin=0, ymin=0, legend pos=north east,
    ]
      \addplot[ochre,no marks] table[col sep=comma,x=prey,y=predator] {predator_prey.csv};
      \addlegendentry{orbite}
      \addplot[only marks,mark=*,mark size=2.2pt,wine] coordinates {(20,10)};
      \addlegendentry{équilibre}
    \end{axis}
    \end{tikzpicture}
    \caption{Portrait de phase}
  \end{subfigure}
  \caption{Solution du modèle de Lotka--Volterra obtenue avec RK4.}\label{fig-lv}
\end{figure}

La \cref{fig-lv} montre des oscillations régulières, de période voisine de $5{,}5$ unités de temps : les proies varient entre $7{,}9$ et $40{,}6$ et les prédateurs entre $3{,}1$ et $23{,}3$. Dans le plan de phase, la trajectoire est une courbe fermée autour de l'équilibre, comme l'annonce la \cref{prop-invariant}.

Pour contrôler la qualité du calcul, nous évaluons l'intégrale première $I$ de la \cref{prop-invariant} le long de la trajectoire échantillonnée : sa variation maximale par rapport à la valeur initiale ($I\approx-3{,}8133$) est inférieure à $1{,}2\times10^{-5}$, du même ordre que l'arrondi à quatre décimales des valeurs enregistrées. La méthode RK4 conserve donc très bien cet invariant sur $[0,30]$.

\section{Discussion}\label{sec-discussion}

Les expériences confirment les prédictions théoriques sur deux points : les ordres de convergence observés sont ceux annoncés, et le coût pour une précision donnée diminue fortement avec l'ordre. Elles ne remplacent pas une analyse de la stabilité : sur un problème raide, c'est la condition de la \cref{sec-stabilite} qui imposerait un pas très petit, quelle que soit la précision visée.

\chapter*{Conclusion générale}
\addcontentsline{toc}{chapter}{Conclusion générale}
\markboth{Conclusion générale}{Conclusion générale}

Ce mémoire a présenté trois méthodes numériques à un pas pour les équations différentielles ordinaires. Nous avons démontré qu'une méthode consistante d'ordre~$p$, dont la fonction d'incrément est lipschitzienne, converge à l'ordre~$p$, avec une estimation explicite de l'erreur globale. Les expériences sur l'équation test ont confirmé les ordres $1$, $2$ et $4$ et mis en évidence l'écart de coût entre les méthodes ; l'application au modèle de Lotka--Volterra a montré que RK4 conserve l'intégrale première du système avec une grande précision.

\paragraph{Perspectives.}
Plusieurs prolongements naturels se dégagent :
\begin{itemize}[itemsep=2pt]
  \item les méthodes à pas adaptatif, qui estiment l'erreur locale pour ajuster $h$ automatiquement ;
  \item les méthodes implicites, mieux adaptées aux problèmes raides, dont la stabilité absolue n'est pas limitée comme au \cref{sec-stabilite} ;
  \item les méthodes géométriques (symplectiques), qui conservent exactement certaines quantités comme l'énergie pour les systèmes hamiltoniens.
\end{itemize}


\appendix
\chapter{Tableaux de Butcher}\label[appendix]{app-butcher}

Une méthode de Runge--Kutta à $s$ étages s'écrit $y_{n+1}=y_n+h\sum_{i=1}^{s}b_ik_i$ avec $k_i=f\bigl(t_n+c_ih,\;y_n+h\sum_{j<i}a_{ij}k_j\bigr)$. Ses coefficients $c_i$, $a_{ij}$ et $b_i$ se présentent dans un \emph{tableau de Butcher} : la première colonne contient les $c_i$, le bloc central la matrice $A=(a_{ij})$ et la dernière ligne les poids $b_i$.

\bigskip
\begin{center}
\renewcommand{\arraystretch}{1.25}
\begin{tabular}{c@{\hspace{3cm}}c}
$\begin{array}{c|cc}
  0 &     &     \\
  1 & 1   &     \\ \hline
    & \tfrac12 & \tfrac12
\end{array}$
&
$\begin{array}{c|cccc}
  0        &          &          &   &   \\
  \tfrac12 & \tfrac12 &          &   &   \\
  \tfrac12 & 0        & \tfrac12 &   &   \\
  1        & 0        & 0        & 1 &   \\ \hline
           & \tfrac16 & \tfrac13 & \tfrac13 & \tfrac16
\end{array}$\\[4pt]
\small Méthode de Heun & \small Méthode RK4 classique
\end{tabular}
\end{center}

\chapter{Exemple de code : la méthode RK4}\label[appendix]{app-code}

La fonction suivante intègre un système $y'=f(t,y)$ avec la méthode de Runge--Kutta d'ordre~4 ; elle correspond à la définition du \cref{sec-classiques}.

\begin{lstlisting}[caption={Méthode RK4 en Python (pas constant).},label={lst-rk4}]
def rk4(f, t0, y0, T, N):
    """Integre y' = f(t, y) sur [t0, T] avec N pas de taille h."""
    h = (T - t0) / N
    t, y = t0, y0
    ys = [y0]
    for _ in range(N):
        k1 = f(t, y)
        k2 = f(t + h / 2, y + h / 2 * k1)
        k3 = f(t + h / 2, y + h / 2 * k2)
        k4 = f(t + h, y + h * k3)
        y = y + h / 6 * (k1 + 2 * k2 + 2 * k3 + k4)
        t = t + h
        ys.append(y)
    return ys
\end{lstlisting}


\backmatter
\printbibliography[heading=bibintoc,title={Bibliographie}]

\end{document}
